\chapter{2001年上海卷（秋文）}
\section{填空题}
\begin{problem}
设函数 \(f(x)=\log_{9}x\)，则满足 \(f(x)=\frac{1}{2}\) 的 \(x\) 值为\fillinblank{}.
\end{problem}

\begin{answer}
\(3\)
\end{answer}

\begin{solution}
由 \(\log_{9}x=\dfrac{1}{2}\)，得
\[
x=9^{1/2}=3
\]
因为对数的真数必须为正，且 \(3>0\)，所以所求 \(x\) 为 \(3\).
\end{solution}

\begin{problem}
设数列 \(\{a_n\}\) 的首项 \(a_1=-7\)，且满足 \(a_{n+1}=a_n+2\)（\(n\in\N\)），则 \(a_1+a_2+\cdots+a_{17}=\fillinblank{}\).
\end{problem}

\begin{answer}
\(153\)
\end{answer}

\begin{solution}
由递推关系可知 \(\{a_n\}\) 是首项为 \(-7\)、公差为 \(2\) 的等差数列，因此
\[
a_{17}=a_1+16\times2=-7+32=25
\]
于是
\[
a_1+a_2+\cdots+a_{17}=\frac{17(a_1+a_{17})}{2}=\frac{17(-7+25)}{2}=153
\]
所以填 \(153\).
\end{solution}

\begin{problem}
设 \(P\) 为双曲线 \(\frac{x^2}{4}-y^2=1\) 上一动点，\(O\) 为坐标原点，\(M\) 为线段 \(OP\) 的中点，则点 \(M\) 的轨迹方程为\fillinblank{}.
\end{problem}

\begin{answer}
\(x^2-4y^2=1\)
\end{answer}

\begin{solution}
设动点 \(P\) 的坐标为 \((u,v)\)，则中点 \(M\) 的坐标为 \((x,y)=\left(\dfrac{u}{2},\dfrac{v}{2}\right)\)，即 \(u=2x\)，\(v=2y\). 由 \(P\) 在双曲线上，得
\[
\frac{(2x)^2}{4}-(2y)^2=1
\]
化简为
\[
x^2-4y^2=1
\]
反过来，若 \((x,y)\) 满足该方程，则点 \((2x,2y)\) 满足原双曲线方程，且 \((x,y)\) 是原点与该点的中点. 因此轨迹方程为 \(x^2-4y^2=1\).
\end{solution}

\begin{problem}
设集合 \(A=\{x\mid 2\lg x=\lg(8x-15),x\in\R\}\)，\(B=\{x\mid \cos\frac{x}{2}>0,x\in\R\}\)，则 \(A\cap B\) 的元素个数为\fillinblank{}\(\text{个}\).
\end{problem}

\begin{answer}
\(1\)
\end{answer}

\begin{solution}
由对数真数的限制，必须有 \(x>0\) 且 \(8x-15>0\)，即 \(x>\dfrac{15}{8}\). 在此条件下，
\[
2\lg x=\lg x^2
\]
所以
\[
x^2=8x-15\Longleftrightarrow (x-3)(x-5)=0
\]
两个根 \(3,5\) 均满足定义域条件，故 \(A=\{3,5\}\).

又因为 \(B=\left\{x\mid\cos\dfrac{x}{2}>0\right\}\)，有
\(\cos\frac{3}{2}>0\)，\(\cos\frac{5}{2}<0\)
其中前者是因为 \(0<\dfrac32<\dfrac\pi2\)，后者是因为 \(\dfrac\pi2<\dfrac52<\pi\). 因而 \(A\cap B=\{3\}\)，元素个数为 \(1\).
\end{solution}

\begin{problem}
抛物线 \(x^2-4y-3=0\) 的焦点坐标为\fillinblank{}.
\end{problem}

\begin{answer}
\(\left(0,\frac{1}{4}\right)\)
\end{answer}

\begin{solution}
将方程整理为
\[
x^2=4\left(y+\frac{3}{4}\right)
\]
与抛物线标准方程 \(x^2=4p(y-y_0)\) 比较，得顶点为 \(\left(0,-\dfrac34\right)\)，且 \(p=1\). 因此焦点为
\[
\left(0,-\frac34+1\right)=\left(0,\frac14\right)
\]
\end{solution}

\begin{problem}
设数列 \(\{a_n\}\) 是公比 \(q>0\) 的等比数列，\(S_n\) 是它的前 \(n\) 项和. \(\displaystyle\lim_{n\to\infty}S_n=7\)，则此数列的首项 \(a_1\) 的取值范围是\fillinblank{}.
\end{problem}

\begin{answer}
\((0,7)\)
\end{answer}

\begin{solution}
若等比数列的前 \(n\) 项和有有限极限，则必须有 \(0<q<1\). 此时
\(S_n=\frac{a_1(1-q^n)}{1-q}\)，\(\lim_{n\to\infty}S_n=\frac{a_1}{1-q}=7\)
从而 \(a_1=7(1-q)\). 当 \(0<q<1\) 时，\(0<a_1<7\).

反过来，任取 \(a_1\in(0,7)\)，令 \(q=1-\dfrac{a_1}{7}\)，则 \(0<q<1\)，且 \(\lim S_n=7\). 所以 \(a_1\) 的取值范围为 \((0,7)\).
\end{solution}

\begin{problem}
某餐厅供应客饭，每位顾客可以在餐厅提供的菜肴中任选 \(2\) 荤 \(2\) 素共 \(4\) 种不同的品种. 现在餐厅准备了 \(5\) 种不同的荤菜，若要保证每位顾客有 \(200\) 种以上不同的选择，则餐厅至少还需要准备不同的素菜品种\fillinblank{}\(\text{种}\)（结果用数值表示）.
\end{problem}

\begin{answer}
\(7\)
\end{answer}

\begin{solution}
设素菜有 \(n\) 种，则每位顾客的选择数为
\[
\mathrm C_5^2\mathrm C_n^2=10\cdot\frac{n(n-1)}{2}=5n(n-1)
\]
题意要求 \(5n(n-1)>200\)，即 \(n(n-1)>40\). 当 \(n=6\) 时，选择数为 \(5\times6\times5=150\)，不超过 \(200\)；当 \(n=7\) 时，选择数为 \(5\times7\times6=210>200\). 因此至少需要准备 \(7\) 种素菜.
\end{solution}

\begin{problem}
在代数式 \(\left(x-\frac{1}{x^2}\right)^5\) 的展开式中，常数项为\fillinblank{}.
\end{problem}

\begin{answer}
\(0\)
\end{answer}

\begin{solution}
展开式的通项为
\(\mathrm C_5^k x^{5-k}\left(-\frac{1}{x^2}\right)^k=(-1)^k\mathrm C_5^k x^{5-3k}\)，\(k=0\)，\(1\)，\(2\)，\(3\)，\(4\)，\(5\).
当 \(k=0\)，\(1\)，\(2\)，\(3\)，\(4\)，\(5\) 时，幂指数依次为 \(5\)，\(2\)，\(-1\)，\(-4\)，\(-7\)，\(-10\)，其中没有 \(0\)，所以展开式中没有常数项，常数项为 \(0\).
\end{solution}

\begin{problem}
设 \(x=\sin\alpha\)，\(\alpha\in\left[-\frac{\pi}{6},\frac{5\pi}{6}\right]\)，则 \(\arccos x\) 的取值范围为\fillinblank{}.
\end{problem}

\begin{answer}
\(\left[0,\frac{2\pi}{3}\right]\)
\end{answer}

\begin{solution}
在区间 \(\left[-\dfrac\pi6,\dfrac{5\pi}{6}\right]\) 内，正弦函数在 \(\alpha=\dfrac\pi2\) 处取得最大值 \(1\)，在左端点处取得最小值 \(-\dfrac12\). 因而
\[
x=\sin\alpha\in\left[-\frac12,1\right]
\]
反余弦函数在 \([-1,1]\) 上单调递减，所以
\[
\arccos x\in\left[\arccos1,\arccos\left(-\frac12\right)\right]
=\left[0,\frac{2\pi}{3}\right]
\]
\end{solution}

\begin{problem}
利用下列盈利表中的数据进行决策，应选择的方案是\fillinblank{}.

\begin{center}
\begin{tabular}{|c|c|c|c|c|c|}
\hline
自然状况 & 概率 & \multicolumn{4}{c|}{盈利（万元）} \\
\cline{3-6}
 & & \(A_1\) & \(A_2\) & \(A_3\) & \(A_4\) \\
\hline
\(S_1\) & \(0.25\) & \(50\) & \(70\) & \(-20\) & \(98\) \\
\hline
\(S_2\) & \(0.30\) & \(65\) & \(26\) & \(52\) & \(82\) \\
\hline
\(S_3\) & \(0.45\) & \(26\) & \(16\) & \(78\) & \(-10\) \\
\hline
\end{tabular}
\end{center}
\end{problem}

\begin{answer}
\(A_3\)
\end{answer}

\begin{solution}
按概率加权计算各方案的期望盈利：
\[
E(A_1)=0.25\times50+0.30\times65+0.45\times26=43.7
\]
\[
E(A_2)=0.25\times70+0.30\times26+0.45\times16=32.5
\]
\[
E(A_3)=0.25\times(-20)+0.30\times52+0.45\times78=45.7
\]
\[
E(A_4)=0.25\times98+0.30\times82+0.45\times(-10)=44.6
\]
其中最大值为 \(E(A_3)=45.7\)，所以应选择方案 \(A_3\).
\end{solution}

\begin{problem}
已知两个圆：
\begin{circlelist}
\item \(x^2+y^2=1\)；
\item \(x^2+(y-3)^2=1\).
\end{circlelist}
由前一个方程减去后一个方程可得上述两圆的对称轴方程. 将上述命题在曲线的情况下加以推广，即要求得到一个更一般的命题，而已知命题应成为所推广命题的一个特例，推广的命题为\fillinblank{}.
\end{problem}

\begin{answer}
若两圆方程为 \((x-a)^2+(y-b)^2=r^2\) 与 \((x-c)^2+(y-d)^2=r^2\)，且 \(a\ne c\) 或 \(b\ne d\)，则由两式相减可得两圆的对称轴方程.
\end{answer}

\begin{solution}
将两个圆的方程分别记为
\((x-a)^2+(y-b)^2=r^2\)，\((x-c)^2+(y-d)^2=r^2\)
其中两个圆心不同. 两式相减，二次项和半径项消去，得到
\[
2(c-a)x+2(d-b)y+a^2+b^2-c^2-d^2=0
\]
这是一条直线，其法向量为 \((c-a,d-b)\)，所以它垂直于两圆心连线；代入圆心连线中点 \(\left(\dfrac{a+c}{2},\dfrac{b+d}{2}\right)\) 可验证该中点在此直线上. 因而这条直线就是两圆的对称轴.

当 \(a=b=0\)，\(c=0\)，\(d=3\)，\(r=1\) 时，上式化为 \(6y-9=0\)，即 \(y=\dfrac32\)，正好得到题设结论. 所以该推广包含原命题.
\end{solution}

\begin{problem}
据报道，我国目前已成为世界上受荒漠化危害最严重的国家之一. 下左图表示我国土地沙化总面积在上个世纪五六十年代，七八十年代，九十年代的变化情况. 由图中的相关信息，可将上述有关年代中，我国年平均土地沙化面积在下右图中图示为：\fillinblank{}

\FigureLayoutDeclare{img/2001/shanghai_liberal/q12_fig1.png}{15.0cm}{22bp 43bp 31bp 71bp}
\bitmapfigure[max width=12cm]{img/2001/shanghai_liberal/q12_fig1.png}
\end{problem}

\begin{answer}
在 \(1950\) 至 \(1970\)、\(1970\) 至 \(1990\)、\(1990\) 至 \(2000\) 年间，纵坐标分别为 \(16\)、\(21\)、\(25\)（百平方公里）.
\end{answer}

\begin{solution}
由左图可知，\(1950\) 至 \(1970\) 年土地沙化总面积增加 \(253.3-250.1=3.2\) 万平方公里，年平均增加 \(3.2\div20=0.16\) 万平方公里，即 \(16\) 百平方公里；\(1970\) 至 \(1990\) 年年平均增加 \((257.5-253.3)\div20=0.21\) 万平方公里，即 \(21\) 百平方公里；\(1990\) 至 \(2000\) 年年平均增加 \((260-257.5)\div10=0.25\) 万平方公里，即 \(25\) 百平方公里. 因此右图应在这三个时间段分别画出高度为 \(16\)、\(21\)、\(25\) 的阶梯图.
\end{solution}

\section{单选题}
\begin{problem}
\(a=3\) 是直线 \(ax+2y+3a=0\) 和直线 \(3x+(a-1)y=a-7\) 平行且不重合的（\quad）

\choices
    {充分非必要条件}
    {必要非充分条件}
    {充要条件}
    {既非充分也非必要条件}
\end{problem}

\begin{answer}
C
\end{answer}

\begin{solution}
若 \(a=3\)，两直线方程分别为
\(3x+2y+9=0\)，\(3x+2y+4=0\).
它们的斜率相同，且常数项不同，因此平行且不重合，故 \(a=3\) 是该命题的充分条件.

反过来，两直线平行的充要条件为
\[
a(a-1)-3\times2=0
\]
即
\[
a^2-a-6=0\Longleftrightarrow (a-3)(a+2)=0
\]
当 \(a=-2\) 时，两式都化为直线 \(-x+y-3=0\)，两直线重合，不符合“不重合”的条件；所以平行且不重合只能有 \(a=3\). 因而 \(a=3\) 也是必要条件，选 C.
\end{solution}

\begin{problem}
在平行六面体 \(ABCD-A_1B_1C_1D_1\) 中，\(M\) 为 \(AC\) 与 \(BD\) 的交点，若 \(\overrightarrow{A_1B_1}=\bs{a}\)，\(\overrightarrow{A_1D_1}=\bs{b}\)，\(\overrightarrow{A_1A}=\bs{c}\)，则下列向量中与 \(\overrightarrow{B_1M}\) 相等的向量是（\quad）

\choices
    {\(-\frac{1}{2}\bs{a}+\frac{1}{2}\bs{b}+\bs{c}\)}
    {\(\frac{1}{2}\bs{a}+\frac{1}{2}\bs{b}+\bs{c}\)}
    {\(\frac{1}{2}\bs{a}-\frac{1}{2}\bs{b}+\bs{c}\)}
    {\(-\frac{1}{2}\bs{a}-\frac{1}{2}\bs{b}+\bs{c}\)}
\end{problem}

\begin{answer}
A
\end{answer}

\begin{solution}
以 \(A_1\) 为原点，用 \(\bs a\)，\(\bs b\)，\(\bs c\) 分别表示 \(\overrightarrow{A_1B_1}\)，\(\overrightarrow{A_1D_1}\)，\(\overrightarrow{A_1A}\). 则
\(\overrightarrow{A_1B_1}=\bs a\)，\(\overrightarrow{A_1A}=\bs c\)，\(\overrightarrow{A_1C}=\bs a+\bs b+\bs c\).
因为平行四边形 \(ABCD\) 的对角线互相平分，\(M\) 是 \(AC\) 的中点，所以
\[
\overrightarrow{A_1M}=\frac12\left(\bs c+\bs a+\bs b+\bs c\right)=\frac12\bs a+\frac12\bs b+\bs c
\]
于是
\[
\overrightarrow{B_1M}=\overrightarrow{A_1M}-\overrightarrow{A_1B_1}
=-\frac12\bs a+\frac12\bs b+\bs c
\]
对应选项 A.
\end{solution}

\begin{problem}
已知 \(a\)，\(b\) 为两条不同的直线，\(\alpha\)，\(\beta\) 为两个不同的平面，且 \(a\perp\alpha\)，\(b\perp\beta\)，则下列命题中的假命题是（\quad）

\choices
    {若 \(a\parallel b\)，则 \(\alpha\parallel\beta\)}
    {若 \(\alpha\perp\beta\)，则 \(a\perp b\)}
    {若 \(a\)，\(b\) 相交，则 \(\alpha\)，\(\beta\) 相交}
    {若 \(\alpha\)，\(\beta\) 相交，则 \(a\)，\(b\) 相交}
\end{problem}

\begin{answer}
D
\end{answer}

\begin{solution}
由于直线垂直于平面时，直线方向是该平面的法向方向：若 \(a\parallel b\)，则两个平面的法向方向平行，结合 \(\alpha\ne\beta\)，可得 \(\alpha\parallel\beta\)，所以 A 正确；若 \(\alpha\perp\beta\)，则两平面的法向方向垂直，所以 \(a\perp b\)，B 正确.

若 \(a\)，\(b\) 相交而 \(\alpha\)，\(\beta\) 不相交，则两个不同平面必平行，从而其法向方向平行，导致 \(a\parallel b\)，与两条不同直线相交矛盾，因此 C 正确.

D 不一定成立. 例如取 \(\alpha:z=0\)，\(\beta:y=0\)，两平面相交；取 \(a\) 为过 \((0,0,1)\) 且方向为 \((0,0,1)\) 的直线，取 \(b\) 为过 \((1,0,0)\) 且方向为 \((0,1,0)\) 的直线. 则 \(a\perp\alpha\)，\(b\perp\beta\)，但 \(a\)，\(b\) 不相交. 所以假命题是 D.
\end{solution}

\begin{problem}
用计算器验算函数 \(y=\frac{\lg x}{x}\)（\(x>1\)）的若干个值，可以猜想下列命题中的真命题只能是（\quad）

\choices
    {\(y=\frac{\lg x}{x}\) 在 \((1,+\infty)\) 上是单调减函数}
    {\(y=\frac{\lg x}{x}\)，\(x\in(1,+\infty)\) 的值域为 \(\left(0,\frac{\lg 3}{3}\right]\)}
    {\(y=\frac{\lg x}{x}\)，\(x\in(1,+\infty)\) 有最小值}
    {\(\displaystyle\lim_{n\to\infty}\frac{\lg n}{n}=0\)，\(n\in\N\)}
\end{problem}

\begin{answer}
D
\end{answer}

\begin{solution}
令 \(f(x)=\dfrac{\lg x}{x}\). 先比较几个定义域内的值：
\(f(2)=\frac{\lg2}{2}\approx0.1505\)，\(f(3)=\frac{\lg3}{3}\approx0.1590\)
所以函数在 \((1,+\infty)\) 上不可能单调递减，选项 A 错. 又
\[
f(2.7)=\frac{\lg2.7}{2.7}\approx0.15976>\frac{\lg3}{3}
\]
故选项 B 所给的值域上界也不正确. 对任意 \(x>1\)，有 \(f(x)>0\)，而 \(\displaystyle\lim_{x\to+\infty}f(x)=0\)，所以函数没有最小值，选项 C 错.

利用对数函数增长速度低于一次函数，
\[
\lim_{x\to+\infty}\frac{\lg x}{x}=0
\]
取 \(x=n\in\N\) 即得
\[
\lim_{n\to\infty}\frac{\lg n}{n}=0
\]
因此只有选项 D 是真命题.
\end{solution}

\section{解答题}
\begin{problem}
已知 \(a\)，\(b\)，\(c\) 是 \(\triangle ABC\) 中 \(\angle A\)，\(\angle B\)，\(\angle C\) 的对边，\(S\) 是 \(\triangle ABC\) 的面积，若 \(a=4\)，\(b=5\)，\(S=5\sqrt{3}\)，求 \(c\) 的长度.
\end{problem}

\begin{solution}
由三角形面积公式，
\[
S=\frac{1}{2}ab\sin C=10\sin C=5\sqrt{3}
\]
所以 \(\sin C=\dfrac{\sqrt{3}}{2}\). 因为 \(0<C<\pi\)，所以 \(C=\dfrac{\pi}{3}\) 或 \(C=\dfrac{2\pi}{3}\). 由余弦定理，
\[
c^2=a^2+b^2-2ab\cos C=41-40\cos C
\]
当 \(C=\dfrac{\pi}{3}\) 时，\(c^2=41-20=21\)，得 \(c=\sqrt{21}\)；当 \(C=\dfrac{2\pi}{3}\) 时，\(c^2=41+20=61\)，得 \(c=\sqrt{61}\). 两个值均满足三角形三边关系，故 \(c=\sqrt{21}\) 或 \(c=\sqrt{61}\).
\end{solution}

\begin{problem}
设 \(F_1\)，\(F_2\) 为椭圆 \(\frac{x^2}{9}+\frac{y^2}{4}=1\) 的两个焦点，\(P\) 为椭圆上的一点. 已知 \(P\)，\(F_1\)，\(F_2\) 是一个直角三角形的三个顶点，且 \(|PF_1|>|PF_2|\)，求 \(\frac{|PF_1|}{|PF_2|}\) 的值.
\end{problem}

\begin{solution}
设 \(r_1=|PF_1|\)，\(r_2=|PF_2|\). 椭圆的长半轴为 \(3\)，由椭圆定义得 \(r_1+r_2=6\)；焦距满足 \(c=\sqrt{3^2-2^2}=\sqrt{5}\)，所以 \(|F_1F_2|=2\sqrt{5}\)，即 \(|F_1F_2|^2=20\).

当直角在 \(P\) 处时，勾股定理给出 \(r_1^2+r_2^2=20\). 结合 \((r_1+r_2)^2=36\)，得 \(r_1r_2=8\)，从而 \(r_1\)，\(r_2\) 是方程 \(t^2-6t+8=0\) 的两个根. 因为 \(r_1>r_2\)，所以 \(r_1=4\)，\(r_2=2\)，此时 \(\dfrac{r_1}{r_2}=2\).

当直角在 \(F_1\) 处时，\(r_2^2=r_1^2+20\). 联立 \(r_1+r_2=6\)，得 \(r_1=\dfrac{4}{3}\)，\(r_2=\dfrac{14}{3}\)，与 \(r_1>r_2\) 矛盾，故此情形不成立.

当直角在 \(F_2\) 处时，\(r_1^2=r_2^2+20\). 联立 \(r_1+r_2=6\)，得 \(r_1=\dfrac{14}{3}\)，\(r_2=\dfrac{4}{3}\)，此时 \(\dfrac{r_1}{r_2}=\dfrac{7}{2}\). 因此所求比值为 \(2\) 或 \(\dfrac{7}{2}\).
\end{solution}

\begin{problem}
在棱长为 \(a\) 的正方体 \(OABC-O'A'B'C'\) 中，\(E\)，\(F\) 分别是棱 \(AB\)，\(BC\) 上的动点，且 \(AE=BF\).
\begin{enumerate}
\item 求证：\(A'F\perp C'E\)；
\item 当三棱锥 \(B'-BEF\) 的体积取得最大值时，求二面角 \(B'-EF-B\) 的大小.（结果用反三角函数表示）
\end{enumerate}
\end{problem}

\begin{solution}
\begin{enumerate}
\item 建立空间直角坐标系，取 \(O=(0,0,0)\)，\(A=(a,0,0)\)，\(B=(a,a,0)\)，\(C=(0,a,0)\)，并取竖直方向为第三坐标轴. 设 \(AE=BF=t\)，其中 \(0\leqslant t\leqslant a\)，则
\(E=(a,t,0)\)，\(F=(a-t,a,0)\)，\(A'=(a,0,a)\)，\(C'=(0,a,a)\).
于是
\(\overrightarrow{A'F}=(-t,a,-a)\)，\(\overrightarrow{C'E}=(a,t-a,-a)\)
从而
\[
\overrightarrow{A'F}\cdot\overrightarrow{C'E}=-at+a(t-a)+a^2=0
\]
所以 \(A'F\perp C'E\).
\item 设 \(BF=AE=x\). 因为 \(EB=a-x\)，且 \(AB\perp BC\)，所以
\[
S_{\triangle BEF}=\frac{1}{2}x(a-x)
\]
点 \(B'\) 到平面 \(BEF\)（即底面 \(OABC\)）的距离为 \(a\)，故三棱锥体积为
\[
V=\frac{1}{3}S_{\triangle BEF}\cdot a=\frac{a}{6}x(a-x)
\]
由 \(x(a-x)\leqslant\dfrac{a^2}{4}\)，且等号当且仅当 \(x=\dfrac{a}{2}\)，可知体积最大时 \(E\)，\(F\) 分别为 \(AB\)，\(BC\) 的中点.

此时 \(BE=BF=\dfrac{a}{2}\)，在直角三角形 \(BEF\) 中，设从 \(B\) 向 \(EF\) 作垂线，垂足为 \(H\). 则
\[
BH=\frac{BE\cdot BF}{EF}=\frac{(a/2)^2}{a/\sqrt{2}}=\frac{a}{2\sqrt{2}}
\]
过 \(B'B\) 和 \(BH\) 的平面垂直于 \(EF\)，其与两个侧面的交线分别为 \(B'H\) 和 \(BH\)，所以二面角 \(B'-EF-B\) 的平面角为 \(\angle B'HB\). 在直角三角形 \(B'BH\) 中，\(BB'=a\)，因此
\[
\tan\angle B'HB=\frac{BB'}{BH}=2\sqrt{2}
\]
故二面角的大小为 \(\arctan\left(2\sqrt{2}\right)\).
\end{enumerate}
\end{solution}

\begin{problem}
对任意一个非零复数 \(z\)，定义集合 \(M_z=\{\omega\mid \omega=z^{2n-1},n\in\N\}\).
\begin{enumerate}
\item 设 \(a\) 是方程 \(x+\frac{1}{x}=\sqrt{2}\) 的一个根，试用列举法表示集合 \(M_a\). 若在 \(M_a\) 中任取两个数，求其和为零的概率 \(P\)；
\item 设集合 \(M_z\) 中只有 3 个元素，试写出满足条件的一个 \(z\) 的值，并说明理由.
\end{enumerate}
\end{problem}

\begin{solution}
\begin{enumerate}
\item 由
\[
x+\frac{1}{x}=\sqrt{2}
\]
得
\[
x^2-\sqrt{2}x+1=0
\]
即
\[
\left(x-\frac{\sqrt{2}}{2}\right)^2+\frac{1}{2}=0
\]
所以方程的两个复数根为
\[
a=\frac{1+\i}{\sqrt{2}}\quad\text{或}\quad a=\frac{1-\i}{\sqrt{2}}
\]
若 \(a=(1+\i)/\sqrt{2}\)，则 \(a^2=\i\)；若 \(a=(1-\i)/\sqrt{2}\)，则 \(a^2=-\i\). 两种情形的奇次幂都得到同一组四个数，因而
\[
M_a=\left\{\frac{1+\i}{\sqrt{2}},\frac{-1+\i}{\sqrt{2}},\frac{-1-\i}{\sqrt{2}},\frac{1-\i}{\sqrt{2}}\right\}
\]
从这四个不同的数中任取两个，所有等可能的取法共有
\(\mathrm C_4^2=6\) 种. 和为零的取法是互为相反数的两对：
\(\left\{\frac{1+\i}{\sqrt{2}},\frac{-1-\i}{\sqrt{2}}\right\}\)，\(\left\{\frac{-1+\i}{\sqrt{2}},\frac{1-\i}{\sqrt{2}}\right\}\).
因此
\[
P=\frac{2}{\mathrm C_4^2}=\frac{1}{3}
\]
\item 取
\[
z=\frac{1}{2}+\frac{\sqrt{3}}{2}\i
\]
则 \(z^2=-\dfrac{1}{2}+\dfrac{\sqrt{3}}{2}\i\)，且 \((z^2)^3=1\)，而 \(z^2\ne1\). 于是 \(z^{2n-1}=z(z^2)^{n-1}\) 按三个值循环，具体为
\[
M_z=\left\{\frac{1}{2}+\frac{\sqrt{3}}{2}\i,-1,\frac{1}{2}-\frac{\sqrt{3}}{2}\i\right\}
\]
这三个数互不相同，所以 \(M_z\) 中恰有 \(3\) 个元素.
\end{enumerate}
\end{solution}

\begin{problem}
用水清洗一堆蔬菜上残留的农药，对用一定量的水清洗一次的效果作如下假定：用 \(1\) 个单位量的水可洗掉蔬菜上残留农药用量的 \(\frac{1}{2}\)，用水越多，洗掉的农药量也越多，但总还有农药残留在蔬菜上. 设用 \(x\) 单位量的水清洗一次以后，蔬菜上残留的农药与本次清洗前残留的农药量之比为函数 \(f(x)\).
\begin{enumerate}
\item 试规定 \(f(0)\) 的值，并解释其实际意义；
\item 试根据假定写出函数 \(f(x)\) 应该满足的条件和具有的性质；
\item 设 \(f(x)=\frac{1}{1+x^2}\)，现有 \(a\)（\(a>0\)）单位量的水，可以清洗一次，也可以把水平均分成 \(2\) 份后清洗两次. 试问用哪种方案清洗后蔬菜上的农药量比较少？说明理由.
\end{enumerate}
\end{problem}

\begin{solution}
\begin{enumerate}
\item 规定 \(f(0)=1\). 这表示使用 \(0\) 单位量的水时没有发生清洗，清洗后农药残留量与清洗前相同.
\item 函数的定义域为 \([0,+\infty)\)，并且 \(f(0)=1\)，\(f(1)=\dfrac{1}{2}\). 由于用水越多，洗掉的农药越多，残留比例应在 \([0,+\infty)\) 上单调递减；又因为清洗后仍有农药残留，所以对任意 \(x\geqslant0\)，有 \(f(x)>0\). 因而 \(0<f(x)\leqslant1\).
\item 设清洗前农药量为 \(A>0\). 一次使用全部水量时，清洗后残留量为
\[
A f(a)=\frac{A}{1+a^2}
\]
将水平均分成两份清洗两次时，每次用水 \(a/2\)，第一次清洗后残留量为 \(A f(a/2)\)，第二次清洗后残留量为
\[
A f\left(\frac{a}{2}\right)^2=\frac{A}{\left(1+\frac{a^2}{4}\right)^2}
\]
因为 \(A>0\)，两次清洗效果更好当且仅当
\[
\frac{1}{\left(1+\frac{a^2}{4}\right)^2}<\frac{1}{1+a^2}
\Longleftrightarrow \left(1+\frac{a^2}{4}\right)^2>1+a^2
\Longleftrightarrow a^2(a^2-8)>0
\]
结合 \(a>0\)，得当 \(a>2\sqrt{2}\) 时两次清洗后残留较少；当 \(a=2\sqrt{2}\) 时两种方案效果相同；当 \(0<a<2\sqrt{2}\) 时一次清洗后残留较少.
\end{enumerate}
\end{solution}

\begin{problem}
对任意函数 \(f(x)\)，\(x\in D\)，可按图示构造一个数列发生器，其工作原理如下：
\begin{circlelist}
\item 输入数据 \(x_0\in D\)，经数列发生器输出 \(x_1=f(x_0)\)；
\item 若 \(x_1\notin D\)，则数列发生器结束工作；若 \(x_1\in D\)，则将 \(x_1\) 反馈回输入端，再输出 \(x_2=f(x_1)\)，并依此规律继续下去.
\end{circlelist}
现定义 \(f(x)=\frac{4x-2}{x+1}\).

\begin{enumerate}
\item 若输入 \(x_0=\frac{49}{65}\)，则由数列发生器产生数列 \(\{x_n\}\). 请写出数列 \(\{x_n\}\) 的所有项；
\item 若要数列发生器产生一个无穷的常数数列，试求输出的初始数据 \(x_0\) 的值；
\item 是否存在 \(x_0\)，在输入数据 \(x_0\) 时，该数列发生器产生一个各项均为负数的无穷数列？若存在，求出 \(x_0\) 的值；若不存在，请说明理由.
\end{enumerate}

\bitmapfigure[max width=6.2cm]{img/2001/shanghai_liberal/q22_fig1.png}
\end{problem}

\begin{solution}
\begin{enumerate}
\item 由于 \(f(x)=\dfrac{4x-2}{x+1}\) 的定义域为 \(D=\{x\in\R\mid x\ne-1\}\)，依次计算得
\[
x_1=f\left(\frac{49}{65}\right)=\frac{\frac{196}{65}-2}{\frac{49}{65}+1}=\frac{11}{19}
\]
\(x_2=f\left(\frac{11}{19}\right)=\frac{\frac{44}{19}-2}{\frac{11}{19}+1}=\frac{1}{5}\)，\(x_3=f\left(\frac{1}{5}\right)=\frac{\frac{4}{5}-2}{\frac{1}{5}+1}=-1\).
因为 \(x_3=-1\notin D\)，发生器在输出 \(x_3\) 后结束，所以数列的所有项为 \(\dfrac{49}{65}\)，\(\dfrac{11}{19}\)，\(\dfrac{1}{5}\)，\(-1\).
\item 要得到无穷的常数数列，初始数据必须是函数的不动点，即 \(x_0=f(x_0)\). 因为 \(x_0\ne-1\)，所以
\[
x_0=\frac{4x_0-2}{x_0+1}\Longleftrightarrow x_0^2-3x_0+2=0\Longleftrightarrow (x_0-1)(x_0-2)=0
\]
因此 \(x_0=1\) 或 \(x_0=2\)，这两个值均在定义域内，且确实分别产生恒为 \(1\)、恒为 \(2\) 的无穷常数数列.
\item 不存在这样的 \(x_0\). 事实上，若某一项 \(x_n<0\)，分两种情况讨论：当 \(x_n<-1\) 时，分子和分母均为负数，故 \(x_{n+1}=f(x_n)>0\)；当 \(-1<x_n<0\) 时，分母为正数，且
\[
f(x_n)<-1\Longleftrightarrow 4x_n-2<-x_n-1\Longleftrightarrow 5x_n<1
\]
所以 \(x_{n+1}<-1\)，再由前一种情况得 \(x_{n+2}>0\). 因此负数项之后至多再出现一项负数，不可能形成各项均为负数的无穷数列.
\end{enumerate}
\end{solution}
